\chapter{Theoretical Background}

\section{Introduction to Project (RDBMS \& ACID)}
A Relational Database Management System (RDBMS) represents the foundational data layer of contemporary enterprise software \cite{silberschatz2020}. Grounded in the mathematical principles of relational calculus and relational algebra formulated by E.~F. Codd \cite{codd1970}, an RDBMS organizes information into formal two-dimensional relations comprising tuples and attributes. Every relation is governed by a formal schema, attribute domain constraints, and primary keys uniquely identifying each entity tuple.

To maintain absolute data integrity amidst concurrent transactions and potential node failures, relational databases uphold the four ACID properties:
\begin{itemize}[itemsep=1.5pt, topsep=2pt, parsep=0pt]
    \item \textbf{Atomicity:} Trip creation and pending invoice generation execute as an indivisible unit of work; failure of either component triggers complete transaction rollback.
    \item \textbf{Consistency:} Declarative foreign key constraints and domain assertions ensure the database transitions exclusively across legally validated states.
    \item \textbf{Isolation:} PostgreSQL Multi-Version Concurrency Control (MVCC) prevents race conditions, dirty reads, and phantom reads during concurrent driver allocation.
    \item \textbf{Durability:} Committed state transitions persist reliably in non-volatile storage via Write-Ahead Logging (WAL) and synchronous disk commits.
\end{itemize}

\section{Existing System Design and Limitations}
Contemporary mobility reservation architectures and naive database implementations suffer from severe structural and operational weaknesses:
\begin{itemize}[itemsep=1.5pt, topsep=2pt, parsep=0pt]
    \item \textbf{Client-Side Calculation Vulnerabilities:} Pricing computed in client-side JavaScript is vulnerable to payload tampering, allowing malicious agents to forge trip fares before submission.
    \item \textbf{Referential Integrity Anomalies:} Soft references devoid of database-enforced foreign keys cause orphaned payment records when reservations are cancelled.
    \item \textbf{Data Redundancy:} Duplicating driver and vehicle details across trip records leads to severe update anomalies and storage bloat.
    \item \textbf{Lack of Cardinality Enforcement:} Enforcing $1:1$ vehicle ownership solely in application logic introduces write race conditions.
    \item \textbf{Inefficient Multi-Table Query Patterns:} Repetitive application round trips cause the $N+1$ query problem under concurrent traffic loads.
\end{itemize}

\section{Proposed Relational Solution}
To resolve these architectural flaws, \textbf{RideShare} implements a resilient database layer:
\begin{itemize}[itemsep=1.5pt, topsep=2pt, parsep=0pt]
    \item \textbf{In-Database Fare Triggers:} Bypasses client computation by embedding distance calculations into a server-side PL/pgSQL \texttt{BEFORE} trigger.
    \item \textbf{Declarative Cardinality:} Enforces strict $1:1$ driver-vehicle ownership through relational \texttt{UNIQUE} and \texttt{FOREIGN KEY} constraints.
    \item \textbf{Cascading Referential Integrity:} Automates invoice cleanup via \texttt{ON DELETE CASCADE} constraints.
    \item \textbf{Optimized Reporting Views:} Provides a unified \texttt{ride\_full\_details} view to eliminate redundant multi-table queries.
\end{itemize}
